\section{Factorization of terminal readers via naturality, refinement, and passage to the limit}

The material starting point is the sole simultaneous constructor that produces a single enriched continuum with five correlated intrinsic operations. The notations \(\mathfrak G_{\rm sp}\), \(\mathfrak G_{\rm sol}\), \(\mathfrak G_{\rm gau}\), \(\mathfrak G_{\rm coh}\), and \(\mathfrak G_{\rm det}\) designate five complete views of that same ledger, reproduced with its fibres, relations, numbers, orientation, carry, memory, boundary, route, multisection, survival, terminals, and readers. They are not taken as five chosen branches or five independent data: their compatibility over the same ledger, their naturality under refinement, their joint non-emptiness, and their passage to the limit are proved, and they are then applied, respectively, to the five problems. No proof in this volume refers to another external source to complete a definition, an identity, or a logical step.

Each closure is presented in a self-contained form: the definitions, proofs, and
compatibilities appear before its classical conclusion. In each resolution,
four distinct layers are exhibited:
\begin{enumerate}
  \item exact identity at every finite level;
  \item compatibility of refinements and preservation of the enriched history;
  \item uniformity and passage to the limit, with its terminal preserved;
  \item recognition of the classical conclusion in the published codomain.
\end{enumerate}

The governing composition of the five closures first preserves the unique
correlative action, forms the continuum only upon passage to the limit, and then opens a view for
each application:
\[
\begin{aligned}
 \mathcal E_k&\xrightarrow{\ \mathcal O_k\ }\mathcal Z_k
 \xrightarrow{\ \varprojlim_k\ }
 \bigl(\mathcal C_{\rm cont}^{\rm disc},\mathcal O_\infty\bigr),\\
 \bigl(\mathcal C_{\rm cont}^{\rm disc},\mathcal O_\infty\bigr)
 &\xrightarrow{\ \operatorname{View}_T\ }\mathfrak G_T^{\rm int}
 \xrightarrow{\ \operatorname{App}_{T,d_T}\ }\mathfrak G_T(d_T),\\
 \mathfrak G_T(d_T)&
 \xrightarrow{\widetilde{\mathcal A}_{T,d_T}}\mathfrak M_T
 \xrightarrow{\mathcal P_T}\mathcal C_T,
 \qquad T\in\{\mathrm{sp,sol,gau,coh,det}\}.
\end{aligned}
\]
Here $\mathcal O_k$ acts once on the complete state and jointly preserves the five features. $\operatorname{View}_T$ neither restricts the domain nor selects histories: it only exposes the coordinates of one application within the complete value of $\mathcal O_\infty$. The classical datum $d_T$ enters only in $\operatorname{App}_{T,d_T}$, after the continuum, and does not participate in generating the intrinsic feature. The real line is likewise a subsequent output of $\operatorname{ev}_{\mathbb R}$, not a premise of the chain.

Each chapter distinguishes the starting domain, the construction, the certificate, the proved conclusion and its exact falsifiers. The preceding source developments construct the five compositions in their respective domains; the following chapters unfold those constructions and carry out the final recognition. No arrow is replaced by a cross-reference, a nominal conclusion or a classical hypothesis.
